\chapter{Refinement of CRE Enforcement}
\label{chapter:refinement-of-cre-enforcement}

%%%%%%%%%%%%%%%%%%%%%%%%%%%%%%%%%%%%%%%%%%%%%%%%%%%%%%%%%%%%%%%%%%%%%%%%%%%%%%%%
%                                                                              %
%%%%%%%%%%%%%%%%%%%%%%%%%%%%%%%%%%%%%%%%%%%%%%%%%%%%%%%%%%%%%%%%%%%%%%%%%%%%%%%%
\section{Introduction}
\label{section:cre-refinement-introduction}
In this chapter we will introduce a different definition of traces (system traces)
and show that, under some conditions, it is isomorph to the (MFOTL) trace
definition introduced in chapter \ref{chapter:preliminaries}.
We will also define what \emph{Refinement} for MFOTL formulae as well as
Lex specifications.
We will show that with our definitions, if a Lex specification A is a refinement
of a Lex specification B, then A compiled to MFOTL is
an MFOTL refinement of B compiled to MFOTL.
We will further introduce a model of runtime enforcement as an event system.
This new model will be based on the new trace definition.
We will give definitions for soundness and transparency of runtime
enforcement event systems.
We will show that these definitions are equivalent to those of
chapter \ref{chapter:preliminaries} under the aforementioned isomorphism
between system and MFOTL traces.
We will adapt the definition of refinement from Sprenger et al. \cite{igloo}
slightly.
Specifically, our mediator function will not only rely on an event being reported
at a certain point in time, but also the trace (history) of previous events.

We will then go on to show that, whenever there exists a sound enforcement
event system for some MFOTL formula $\phi$, and another formula $\psi$ refines
$\phi$, then there exists a sound enforcement event system for $\psi$.
Separately we will show that there exists a sound enforcement event system
for some MFOTL formula $\phi$ if and only if there exists a sound 
enforcer, as defined in chapter \ref{chapter:preliminaries} for $\phi$.
Hublet et al. \cite{lex} have shown that an causable Lex specification
compiles to an causable MFOTL formula.
And Hublet et al. \cite{whyenf} have shown that if an MFOTL formula is 
causable, there exists a sound enforcer for it.
Putting everything together we will have shown that if a Lex specification A
refines an enforceable Lex specification B, then B is enforceable.


\section{Updated Enforcement Model}
\label{section:updated-enforcement-model}

%%%%%%%%%%%%%%%%%%%%%%%%%%%%%%%%%%%%%%%%%%%%%%%%%%%%%%%%%%%%%%%%%%%%%%%%%%%%%%%%
\subsection{System Traces}
\label{subsection:system-traces}
We introduce a different trace definition as in Definition \ref{def:trace}
and we will show that this new definition is at least as expressive as
the old one.

\begin{definition}
    \label{def:system-sequence}
    A \emph{System Sequence} $\sigma_\sys \in \SeqSys$ 
    is a finite or infinite sequence of events
    \begin{equation*}
        \Sequence[1 \leq i \leq k]{e_i},
        k \in \Nat \cup \lbrace \infty \rbrace
    \end{equation*}
    where for every $i \in [1, k]$,
    \begin{equation*}
        e_i \in  \Databases \cup \left\{ \tick \right\}.
    \end{equation*}
\end{definition}
The special event $\tick$ represents a clock tick in the system.
As a more compact notation we use
$$e^k := \underbrace{e \cdot \ldots \cdot e}_{k \text{ times}}$$ for all $k \in \Nat$
and $e \in \Databases \cup \{tick\}$. Note that for $k=0$, $e^k = \emptytrace_\sys$.

\begin{definition}
    \label{def:num-ticks}
    The function
    \begin{equation*}
        \numTicks : \SeqSys \to \Nat \cup \{ \infty \}
    \end{equation*}
    \begin{align*}
        \numTicks(\sigma_\sys) =
        \begin{cases}
            % 0 &\text{if } \sigma_\sys = \emptytrace_\sys \\
            k
                &\text{if } \sigma_\sys = \tick^k
                \text{ for } k \geq 0 \\
            k + \numTicks(\sigma'_\sys)
                &\text{if } \sigma_\sys = \tick^k \cdot D \cdot \sigma'_\sys
                \text{ for } k \geq 0
                \text{ and } D \in \Databases
        \end{cases}
    \end{align*}
    gives the number of ticks in a system sequence.
\end{definition}

\begin{definition}
    \label{def:trim-system-traces}
    The function
    \begin{equation*}
        \trim: \SeqSys \to \SeqSys
    \end{equation*}
    \begin{align*}
        \trim(\sigma_\sys) :=
        \begin{cases}
            \emptytrace_\sys
                &\text{if } \sigma_\sys = \tick^k
                \text{ for } k \geq 0 \\
            \tick^k \cdot D \cdot \trim(\sigma_\sys')
                &\text{if } \sigma_\sys = \tick^k \cdot D \cdot \sgima_\sys'
                \text{ for } k \geq 0
                \text{ and } D \in \Databases
        \end{cases}
    \end{align*}
    removes all $\tick$ events from the end of a system sequence
    that are not followed by a database.
\end{definition}

For sets of system sequences $A \subseteq \SeqSys$, 
$\trim(A) := \{ \trim(\sigma) \mid \sigma \in A \}$.

\begin{definition}
    \label{def:system-trace}
    A \emph{System Trace}
    $\sigma_\sys \in \Traces_\sys \subseteq \SeqSys$ 
    is a finite or infinite system sequence where
    the number of ticks progresses with the number of databases
    (if $\numDBs(\sigma_\sys) = \infty$, then
    $\numTicks(\trim(\sigma_\sys)) = \infty $).
\end{definition}

\begin{definition}
    \label{def:num-db}
    The function
    \begin{equation*}
        \numDBs : \Traces_\sys \to \Nat \cup \{ \infty \}
    \end{equation*}
    \begin{align*}
        \numDBs(\sigma_\sys) =
        \begin{cases}
            0 &\text{if } \sigma_\sys = \tick^k \text{ and } k \geq 0 \\
            1 + \numDBs(\sigma'_\sys)
                &\text{if } \sigma_\sys = \tick^k \cdot D \cdot \sigma'_\sys
                \text{ for } k \geq 0
                \text{ and } D \in \Databases
        \end{cases}
    \end{align*}
    gives the number of database events in a system trace.
\end{definition}


We say a system trace $\sigma_\sys \in \Traces_\sys$
is finite iff $\numDBs(\sigma_\sys) \in \Nat$
and it is infinite iff $\numDBs(\sigma_\sys) = \infty$.
We use
$\Traces_{\sys f} \subseteq 
\left( \Databases \cup \left\{\tick \right\} \right)^*$,
$\Traces_{\sys \omega} \subseteq
\left( \Databases \cup \left\{\tick \right\} \right)^\infty$,
and
$\Traces_\sys = 
\Traces_{\sys f} \cup \Traces_{\sys \omega}$
to denote the sets of finite system traces, infinite system traces,
and all system traces respectively.
Whenever it is clear that the system traces are meant and not the
MFOTL traces as in Section \ref{section:mfotl}, we omit the
subscript "$\sys$".
We also use $\emptytrace_\sys$ to denote an empty system
trace and will omit the subscript if the meaning is clear from
context.



\begin{definition}
    \label{def:cconv'}
    The \emph{Finite Extended System Trace Converter Function} converts system
    traces starting at some timestamp to MFOTL traces.
    \begin{align*}
        \cConv':
        \Nat \times \Traces_\sys
        \to \Traces
    \end{align*}
    \begin{align*}
        % \label{eq:cconv'}        
        \cConv'(\tau, \sigma_\sys) :=
        \begin{cases}
            \emptytrace
                % & \text{if } \trim(\sigma_\sys) = \emptytrace_\sys \\
                & \text{if } \trim(\sigma_\sys) = 0 \\
            (\tau + k, D) \cdot \cConv'(\tau+k, \sigma'_\sys)
                & \text{if } \sigma_\sys = \tick^k \cdot D \cdot \sigma'_\sys \\
                & \text{with } k \geq 0, D \in \Databases \\
        \end{cases}
    \end{align*}
\end{definition}

\begin{definition}
    \label{def:cconv}
    The \emph{System Trace Converter Function} converts system
    traces starting at timestamp $0$ to MFOTL traces.
    \begin{equation*}
        \cConv: \Traces_\sys \to \Traces
    \end{equation*}
    \begin{align*}
        % \label{eq:cconv}
        \cConv(\sigma_\sys) := \cConv'(0, \sigma_\sys).
    \end{align*}
\end{definition}

\begin{definition}
    \label{def:dconv'}
    The extended \emph{MFOTL Trace Converter Function} converts MFOTL
    traces starting at some timestamp to system traces.
    \begin{align*}
        \dConv':
        \Nat \times \Traces \to \trim(\Traces_\sys)
    \end{align*}

    \begin{align*}
        \label{eq:dconv'}
        \dConv'(\tau, \sigma) :=
        \begin{cases}
            \emptytrace_\sys
                &\text{if } \sigma = \emptytrace \\
            \tick^k \cdot D \cdot \dConv'(\tau+k, \sigma')
                &\text{if } \sigma = (\tau+k, D) \cdot \sigma' \\
                &\text{with } k \geq 0
        \end{cases}
    \end{align*}
\end{definition}

\begin{definition}
    \label{def:dconv}
    The \emph{MFOTL Trace Converter Function} converts MFOTL
    traces starting at timestamp $0$ to system traces.
    \begin{equation*}
        \dConv: \Traces \to \trim(\Traces_\sys)
    \end{equation*}
    \begin{equation*}
        \label{eq:dconv}
        \dConv(\sigma) := \dConv'(0, \sigma).
    \end{equation*}
\end{definition}

It is important to note that $\dConv'(\tau, \sigma)$
is only defined, if
$\tau \leq \fts(\sigma)$ or $\sigma = \emptytrace$.
$\dConv$ however, is defined for all $\sigma \in \Traces$,
since timestamps in any trace must increase monotonically,
and the timestamp argument passed to $\dConv'$ is $0$.
To see that the images of both $\dConv$ and $\dConv'$
for all allowed inputs is $\trim(\Traces_\sys)$ and
not $\Traces_\sys$, note that a $\tick$ is only added when
the trace is not yet empty, thus any inserted $\tick$ will
eventually be followed by some database $D$.
For any set of MFOTL traces $A \subseteq \Traces$ and any
set of system traces $B \in \Traces_\sys$ and timestamp
$\tau \in \Nat$, we write
\begin{align*}
    &\cConv(A) &&:= \{ \cConv(\sigma_\sys) \mid \sigma_\sys \in A \}, \\
    &\cConv'(\tau, A) &&:= \{ \cConv'(\tau, \sigma_\sys) \mid \sigma_\sys \in A \}, \\
    &\dConv(B) &&:= \{ \dConv(\sigma) \mid \sigma \in B \}, and \\
    &\dConv'(\tau, B) &&:= \{ \dConv'(\tau, \sigma) \mid \sigma \in B \}. \\
\end{align*}

\begin{theorem}
    \label{theo:trace-equivalence}
    The trace converter functions $\cConv$ and $\dConv$ define a
    bijective mapping between MFOTL traces $\Traces$ and
    trimmed system traces $\trim(\Traces_\sys)$
    and are each others inverses.
    \begin{equation*}
        \forall \sigma \in \Traces .~
        \sigma = \dConv(\cConv(\sigma))
    \end{equation*}
    and
    \begin{equation*}
        \forall \sigma_\sys \in \trim(\Traces_\sys) .~
        \sigma_\sys = \cConv(\dConv(\sigma))
    \end{equation*}
\end{theorem}

The intuition behind the necessity of trimming system traces to make
them fully comparable with MFOTL traces is that MFOTL traces only capture
some last timestamp at which a database was logged, while system traces
can capture a current timestamp (the total number of ticks in a system trace)
which may be greater than the last time stamp at which a database was logged
(the number of ticks up until the last database).

\begin{proof}
    \label{proof:trace-equivalence}
    The proof for Theorem \ref{theo:trace-equivalence}
    proceeds by splitting it into four different properties.
    First we show
    \begin{align}
        \label{eq:domain-finite-mfotl-traces}
        &\forall \sigma \in \Traces_f .~
        &\sigma
        &=\dConv(\cConv(\sigma)),
        \\
        \label{eq:domain-infinite-mfotl-traces}
        &\forall \sigma \in \Traces_\omega .~
        &\sigma
        &=\dConv(\cConv(\sigma)),
        \\
        \label{eq:domain-finite-system-traces}
        &\forall \sigma_\sys \in \trim(\Traces_{\sys f}) .~
        &\sigma_\sys 
        &=\cConv(\dConv(\sigma_\sys)),
        \\
        \label{eq:domain-infinite-system-traces}
        &\forall \sigma_\sys \in \trim(\Traces_{\sys \omega}) .~
        &\sigma_\sys 
        &=\cConv(\dConv(\sigma_\sys)).
    \end{align}
    The proofs for each part actually depend on more general properties
    over the functions $\cConv$ and $\dConv'$ for all $\tau \in \Nat$
    \begin{align}
        \label{eq:domain-finite-mfotl-traces'}
        &\forall \sigma \in \Traces_f .~
        ( \sigma = \emptytrace \vee  \tau \leq \fts(\sigma) )
        \to 
        \cConv'(\tau, \dConv'(\tau, \sigma)) = \sigma,
        \\
        \label{eq:domain-infinite-mfotl-traces'}
        &\forall \sigma \in \Traces_\omega .~
        ( \sigma = \emptytrace \vee  \tau \leq \fts(\sigma) )
        \to 
        \cConv'(\tau, \dConv'(\tau, \sigma)) = \sigma,
        \\
        \label{eq:domain-finite-system-traces'}
        &\forall \sigma_\sys \in \trim(\Traces_{\sys f}) .~
        \dConv'(\tau, \cConv'(\tau, \sigma_\sys)) = \sigma_\sys,
        \\
        \label{eq:domain-infinite-system-traces'}
        &\forall \sigma_\sys \in \trim(\Traces_{\sys \omega}) .~
        \dConv'(\tau, \cConv'(\tau, \sigma_\sys)) = \sigma_\sys.
    \end{align}
    When fixing $\tau=0$ in the four equations above, the preconditions
    all evaluate to "true" and the occurrences of $\cConv'$ and $\dConv'$
    can be replaced with $\cConv$ and $\dConv$ respectively and thus
    Equation \ref{eq:domain-finite-mfotl-traces'} implies
    Equation \ref{eq:domain-finite-mfotl-traces},
    Equation \ref{eq:domain-infinite-mfotl-traces'} implies
    Equation \ref{eq:domain-infinite-mfotl-traces},
    Equation \ref{eq:domain-finite-system-traces'} implies
    Equation \ref{eq:domain-finite-system-traces}, and
    Equation \ref{eq:domain-infinite-system-traces'} implies
    Equation \ref{eq:domain-infinite-system-traces}.

    \begin{enumerate}
        \item \textbf{Finite MFOTL Traces}
            (Equations \ref{eq:domain-finite-mfotl-traces} and
            \ref{eq:domain-finite-mfotl-traces'})
        
        We proof this by structural induction over traces $\sigma \in \Traces_f$
        for all $\tau \in \Nat$ where either $\tau \leq \fts(\sigma)$ or
        $\sigma = \emptytrace$.
    
        \textbf{Base Case:}
        $\sigma = \emptytrace$ (and some $\tau \in \Nat$)
        \begin{align*}
            \cConv'(\tau, \dConv'(\tau, \sigma))
            &= \cConv'(\tau, \dConv(\tau, \emptytrace)) \\
            &= \cConv'(\tau, \emptytrace_\sys) \\
            &= \emptytrace = \sigma
        \end{align*}
        \textbf{Induction Hypothesis:} 
        For every true suffix $\sigma'$ of a trace $\sigma \in \Traces_f$
        that is not the entire trace $\sigma$, i.e.
        $\sigma = \sigma'' \cdot \sigma'$ and $\sigma'' \neq \emptytrace$,
        and for every timestamp $\tau \in \Nat$
        which is less than the first timestamp
        ($\tau \leq \fts(\sigma')$)
        in $\sigma'$, unless $\sigma' = \emptytrace$ it holds that
        \begin{equation*}
            \sigma' = \cConv'(\tau, \dConv'(\tau, \sigma')).
        \end{equation*}
    
        \textbf{Induction Step:}
        Let $\sigma \neq \emptytrace$
        and $\tau \leq \fts(\sigma)$.
        Because $\sigma$ is not empty,
        $\sigma = \left(\tau+k, D \right) \cdot \sigma'$
        for some $k \geq 0$, some $D \in \Databases$ and
        a suffix $\sigma' \in \Traces_f$.
        \begin{align*}
            \cConv'(\tau, \dConv'(\tau, \sigma))
            &= \cConv'(\tau, \dConv'(\tau, (\tau+k, D) \cdot \sigma')) \\
            &\stackrel{\ref{def:dconv'}}{=}
                \cConv'(\tau, \tick^k \cdot D \cdot \dConv'(\tau + k, \sigma')) \\
            &\stackrel{\ref{def:cconv'}}{=}
                (\tau+k, D) \cdot \cConv'(\tau+k, \dConv'(\tau+k, \sigma')) \\
            &\stackrel{\text{IH}}{=} \left(\tau+k, D \right) \cdot \sigma' \\
            &= \sigma
        \end{align*}
    
        We have now shown \ref{eq:domain-finite-mfotl-traces} 
        which does imply \ref{eq:domain-finite-mfotl-traces'}.
        
        \item \textbf{Infinite MFOTL Traces}
            (Equations \ref{eq:domain-infinite-mfotl-traces} and 
            \ref{eq:domain-infinite-mfotl-traces'})
            
        We have already shown that the property holds for arbitrarily
        long finite prefixes of any infinite trace.
        

        \begin{equation*}
            TODO
        \end{equation*}
        
        \item \textbf{Finite System Traces}
            (Equations \ref{eq:domain-finite-system-traces} and 
            \ref{eq:domain-finite-system-traces'})

        This proof proceeds by structural induction over finite and trimmed system traces
        $\sigma_\sys \in \trim(\Traces_{\sys f})$ for all $\tau \in \Nat$
        .
        % where either
        % $\sigma_\sys = \emptytrace_\sys$ or $\tau \leq \fts_\sys(\sigma_\sys)$.
            \begin{enumerate}
                \item \textbf{Base Case} $\sigma_\sys = \emptytrace_\sys$

                \begin{align*}
                    \dConv'(\tau, \cConv'(\tau, \sigma_\sys))
                    &=  \dConv'(\tau, \cConv'(\tau, \emptytrace_\sys)) \\
                    &\stackrel{\ref{def:cconv'}}{=}
                        \dConv'(\tau, \emptytrace) \\
                    &\stackrel{\ref{def:dconv'}}{=}
                        \emptytrace_\sys \\
                    &= \sigma_\sys
                \end{align*}

                \item \textbf{Induction Hypothesis}
                For every true suffix $\sigma_\sys'$ of a trimmed system trace
                $\sigma_\sys \in \trim(\Traces_{\sys f})$
                that is not the entire trace $\sigma_\sys$, i.e.
                $\sigma_\sys = \sigma_\sys'' \cdot \sigma_\sys'$ and
                $\sigma_\sys'' \neq \emptytrace$,
                and for every timestamp $\tau \in \Nat$
                % which is less than the first timestamp
                % ($\tau \leq \fts(\sigma_\sys')$)
                % in $\sigma'$, unless $\sigma' = \emptytrace$
                it holds that
                \begin{equation*}
                    \sigma_\sys'
                    = \dConv'(\tau, \cConv'(\tau, \sigma_\sys')).
                \end{equation*}

                \item \textbf{Induction Step}

                Let $\sigma_\sys \in \trim(\Traces_\sys)$,
                with $\sigma_\sys \neq \emptytrace_\sys$
                .
                % and $\tau \leq \fts_\sys(\sigma_\sys)$.
                Then $\sigma_\sys = \tick^k \cdot D \cdot \sigma_\sys'$
                for some $k \geq 0$, $D \in \Databases$,
                and a suffix $\sigma_\sys' \in \trim(\Traces_\sys)$.
                \begin{align*}
                    \dConv'(\tau, \cConv'(\tau, \sigma_\sys)) &=
                    \dConv'(\tau, \cConv'(\tau, \tick^k \cdot D \cdot \sigma_\sys')) \\
                    &\stackrel{\ref{def:cconv'}}{=}
                        \dConv'(\tau, (\tau+k, D) \cdot \cConv'(\tau+k, \sigma_\sys')) \\
                    &\stackrel{\ref{def:dconv'}}{=}
                        \tick^k \cdot D \cdot
                        \dConv'(\tau+k, \cConv'(\tau+k, \sigma_\sys')) \\
                    &\stackrel{IH}{=}
                        \tick^k \cdot D \cdot \sigma_\sys' \\
                    &= \sigma_\sys
                \end{align*}
                
            \end{enumerate}

        
        \item \textbf{Infinite System Traces}
            (Equations \ref{eq:domain-infinite-system-traces} and 
            \ref{eq:domain-infinite-system-traces'})
            
        \begin{equation*}
            TODO
        \end{equation*}
        
    \end{enumerate}

\end{proof}
Intuitively $\cConv$ and $\cConv'$ count the clock ticks and store them
explicitly as integer timestamps.
We believe that it is clear which definition is used based on the context.
At times we may use the function $\cConv$ to make it more explicit.
For both definitions, a property is a subset of all possible traces.
For conversion of properties we use $\cConv$ with arguments being sets of
traces in which case $\cConv$ simply gets mapped on each element of the set.
As in Sprenger et. al \cite{igloo} and Abadi and Lamport \cite{refinement-mappings}
traces of this second form can be extended with a stuttering step
$\sskip$, where traces are equivalent if they are the same with all
$\sskip$ steps removed.
In particular this means that sets of system traces (such as trace properties)
are closed under the addition and removal of $\sskip$ events.
Thus the set of events in system traces $\Traces_\sys$
gets extended to be
$$
    \Databases \cup \left\{ \tick, \sskip \right\}.
$$
Since $\sskip$ events do not meaningfully modify a system trace,
the isomorphism from Theorem \ref{theo:trace-equivalence} works if
we remove all $\sskip$ events before doing the trace conversion.

With the isomorphism between MFOTL traces as defined in prior work and
Section \ref{chapter:preliminaries} and trimmed system traces we can
define MFOTL semantics for infinite system traces as followed.

\begin{definition}
    \label{def:mfotl-semantics-for-system-traces}
    The semantics of MFOTL for a fixed, infinite system trace
    $\sigma_\sys \in \Traces_{\sys \omega}$, a valuation $v$,
    time-point $i$, and MFOTL formula $\phi$ is defined as
    \begin{equation*}
        v, i \vDash_{\sigma_\sys} \phi \quad \text{ iff } \quad
        v, i \vDash_{\cConv(\sigma_\sys)} \phi.
    \end{equation*}
\end{definition}
We omit the subscript of a trace, if it is clear from context
what trace is meant.

%%%%%%%%%%%%%%%%%%%%%%%%%%%%%%%%%%%%%%%%%%%%%%%%%%%%%%%%%%%%%%%%%%%%%%%%%%%%%%%%
\subsection{Runtime Enforcement as an Event System}
\label{subsection:runtime-enforcement-event-system}


We use a slightly different model of an enforcer as that of
Hublet et al. \cite{whyenf} described in subsection
\ref{subsection:proactive-runtime-enforcement}.
We take a less declarative approach that considers the enforcer as 
a function. And then we use this function to describe an event system.
This yields a similar model as Edit Automata in Ligatti et al.
\cite{edit-automata}.

Our model of proactive enforcement as an event system relies on two 
update functions, one for the proactive and one for the reactive case

\begin{definition}
    \label{def:reactive-update-function}
    A \emph{Reactive Update Function}
    $$
       \RFun : \Traces_{\sys f} \times \Databases \to \Databases
    $$
    modifies a reported  set of events by suppressing
    (deleting) and causing (inserting) certain events
    if needed.
    Where for all $\sigma \in \Traces_{\sys f}$
    and $D \in \Databases$
    \begin{itemize}
        \item $D \setminus \RFun(\sigma, D) \in \Databases_\Sup$
        \item $\RFun(\sigma, D) \setminus D \in \Databases_\Cau$
    \end{itemize}
\end{definition}

\begin{definition}
    \label{def:proactive-update-function}
    A \emph{Proactive Update Function}
    $$
       \PFun :
       \Traces_{\sys f}
       \to
       \Databases_\Cau \cup \left\{ \sskip \right\}
    $$
    indicates whether or not some events must be caused (inserted)
    before a clock tick happens.
\end{definition}

\begin{definition}
    \label{def:enf-state-space}
    The \emph{Enforcement State Space} $\Statespace = \Traces_\sys$
    consists of all system traces (that may include $\sskip$).
\end{definition}

\begin{definition}
    \label{def:enf-initial-states}
    The single initial state is the empty trace $\emptytrace_\sys$.
    $$\Initialstates = \lbrace \emptytrace_\sys \rbrace$$
\end{definition}

\begin{definition}
    \label{def:enf-event-space}
    The \emph{Event Space}
    \begin{equation*}
        \Eventspace=\Databases \cup \lbrace \tick, \sskip \rbrace
    \end{equation*}
    consists of sets of events or the two special events $\tick$ and $\sskip$.
\end{definition}


\begin{definition}
    \label{def:enf-transition-relation}
    The \emph{Transition Relation} is a subset
    \begin{equation*}
        \Transitionrelation{}{} \subseteq
        \Traces_\sys
        \times \left( \Databases \cup \left\{\tick, \sskip \right\} \right)
        \times \Traces_\sys
    \end{equation*}
    The transition relation is by defined  the following three transition
    rules.
    \begin{align*}
        \sigma &\Transitionrelation{}{\sskip} \sigma \cdot \sskip
            & \text{(Skip)} \\
        \sigma &\Transitionrelation{}{\RFun(\sigma, D)} \sigma \cdot \RFun(\sigma, D)
            & \text{(Step-R)} \\
        \sigma &\Transitionrelation{}{\PFun(\sigma) \cdot \tick}
            \sigma \cdot \PFun(\sigma) \cdot \tick
            & \text{(Step-P)}
    \end{align*}
\end{definition}
Note that the "Skip" transition is equivalent to
$\sigma \Transitionrelation{}{\sskip} \sigma$, since system traces
are closed under the addition of $\sskip$ events.
We write $\sigma \Transitionrelation{}{e} \sigma'$
as a shorthand for
$\left(\sigma, e, \sigma'\right) \in \Transitionrelation{}{}$.
Further for events $e_1, e_2 \in \Eventspace$,
$\sigma \Transitionrelation{}{e_1 \cdot e_2} \sigma'$ iff
there exists $\sigma''$ such that
$\sigma \Transitionrelation{}{e_1} \sigma''$ and
$\sigma'' \Transitionrelation{}{e_2} \sigma'$.
This syntax can be extend to sequences of events (i.e. system traces)
$\rho_1, \rho_2 \in \Traces_\sys$,
$\sigma \Transitionrelation{}{\rho_1 \cdot \rho_2} \sigma'$ iff
there exists $\sigma''$ such that
$\sigma \Transitionrelation{}{\rho_1} \sigma''$ and
$\sigma'' \Transitionrelation{}{\rho_2} \sigma'$.

\begin{definition}
    \label{def:enforcement-event-system}
    An \emph{Enforcement Event System} 
    $\EventSystem = \left(\Statespace, \Eventspace, \RFun, \PFun \right)$
    is a tuple consisting of a state space $\Statespace := \Traces_\sys$,
    event space $\Eventspace = \Databases \cup \left\{ \tick, \sskip \right\}$,
    a reactive update function
    $\RFun : \Traces_{\sys f} \times \Databases \to \Databases$
    and a proactive update function
    $\PFun : \Traces_{\sys f} \to \Databases_\Cau \cup \left\{ \sskip \right\}$.
    Where the two update functions define a transition relation as 
    defined in Definition \ref{def:enf-transition-relation}.
\end{definition}

We use the same definition of event traces for an enforcement event system 
$\EventSystem$ with initial states $\Initialstates$ as for
event systems inin Sprenger et al. \cite{igloo} (Section \ref{section:event-systems}),
states $\Initialstates$.
Our described system has the property that if the only initial state is
the empty trace $\emptytrace_\sys$, the final state of the event system
and the event trace of the system are identical.

We define soundness notions for enforcement event systems that are analogous to
soundness of an enforcer in Subsection
\ref{subsection:proactive-runtime-enforcement}.
Sprenger et al. \cite{igloo} defines traces as always being finite sequences of
events. We extend that definition to include both finite and infinite
event sequences.
The function $\EStraces(\EventSystem, \Initialstates)$ then denotes both all
finite and infinite event sequences of the event system $\EventSystem$
starting in any initial state in $\Initialstates$.
$\EStraces_f(\EventSystem, \Initialstates)$ denotes only the finite traces
and $\EStraces_\omega(\EventSystem, \Initialstates)$ only the
infinite traces.
We also extend the definition of trace properties to include both
finite and infinite sequences.
We use the following notations if we specifically talk
about trace properties that are strictly finite or infinite respectively:
$\EventSystem, \Initialstates \vDash_f P$ iff
$\EStraces_f(\EventSystem, \Initialstates) \subseteq P$
and
$\EventSystem, \Initialstates \vDash_\omega P$ iff
$\EStraces_\omega(\EventSystem, \Initialstates) \subseteq P$.
Like in Sprenger et. al \cite{igloo} we will omit the set braces
for initial states, if the set is a singleton.
Note that system traces are only considered infinite if they
are infinite with all skips removed and trimmed (ticks not followed
by a database removed).

Next we will show that any enforcer as in Hublet et al. \cite{whyenf} has a
corresponding enforcement event system.
To simplify this we define a helper function, that can compute the
state of an enforcer after having enforced some trace as well as potentially
some "leftover" ticks.

\begin{definition}
    \label{def:enforcer-state-function}
    The \emph{Enforcer State Function}
    \begin{equation*}
        \state : \EnfStates \times \Traces_\sys \to \EnfStates
    \end{equation*}
    gives the state starting in some initial state, after having enforced
    the corresponding MFOTL trace of a system trace as well as the ticks
    that may or may not be "dangling" at the end of the system trace.
    This function is defined by essentially simulating the $\run$ function of
    Algorithm \ref{alg:loop}.
    The simulation uses a helper function
    \begin{equation*}
        \state' : \EnfStates \times \Traces_\sys \times \Traces_\sys \to \EnfStates
    \end{equation*}
    $\state'(s, \sigma_{\sys 1}, \sigma_{\sys 2}) := $
    \begin{align*}
        \begin{cases}
            s
                &\text{if } \sigma_{\sys 2} = \emptytrace_\sys \\
            \state'(s', \sigma_{\sys 1} \cdot D_\Cau \cdot \tick, \sigma_{\sys 2}')
                &\text{if } \sigma_{\sys 2} = \tick \cdot \sigma_{\sys 2}' \\
                &\text{and } (\PCom(D_\Cau), s') \\
                &\quad = \EnfFun(\cConv(\sigma_{\sys 1}), s, \numTicks(\sigma_{\sys 1}), \tick) \\
            \state'(s', \sigma_{\sys 1} \cdot \tick, \sigma_{\sys 2}')
                &\text{if } \sigma_{\sys 2} = \tick \cdot \sigma_{\sys 2}' \\
                &\text{and } (\NCom, s') \\
                &\quad = \EnfFun(\cConv(\sigma_{\sys 1}), s, \numTicks(\sigma_{\sys 1}), \tick) \\
            \state'(s', \sigma_{\sys 1} \cdot (D \setminus D_\Sup \cup D_\Cau), \sigma_{\sys 2}')
                &\text{if } \sigma_{\sys 2} = D \cdot \sigma_{\sys 2}' \\
                &\text{and } (\RCom(D_\Cau, D_\Sup), s') \\
                &\quad = \EnfFun(\cConv(\sigma_{\sys 1}), s, \numTicks(\sigma_{\sys 1}), D)
        \end{cases}
    \end{align*}
    With that $\state$ is defined as followed for any enforcer state $s \in \EnfStates$
    and system trace~$\sigma_\sys \in \Traces_\sys$
    \begin{equation*}
        \state(s, \sigma_\sys) := \state'(s, \emptytrace_\sys, \sigma_\sys).
    \end{equation*}
\end{definition}

\begin{definition}
    \label{def:induced-event-system}
    The \emph{induced enforcement event system} of an enforcer
    $\Enforcer = (\EnfStates, s_0, \EnfFun)$ is defined as
    $ \EventSystem(\Enforcer) :=
    (\Traces_\sys, \Databases \cup \{\tick, \sskip \}, \RFun, \PFun), $
    with  $\RFun$ and $\PFun$ defined as followed
    \begin{align*}
        \RFun(\sigma_\sys, D)
        &:= D \setminus D_\Sup \cup D_\Cau \\
        &\qquad \text{with }
        \EnfFun(\cConv(\sigma_\sys), \state(s_0, \sigma_\sys), \numTicks(\sigma_\sys), D) =
        (\RCom(D_\Cau, D_\Sup), s) % TODO: how to involve underlying system state??
        \\
        \PFun(\sigma_\sys) &:=
        \begin{cases}
            \sskip &\text{if }
                \EnfFun(\cConv(\sigma_\sys), \state(s_0, \sigma_\sys), \numTicks(\sigma_\sys), D) =
                (\NCom, s) \\
            D_\Cau &\text{if }
                \EnfFun(\cConv(\sigma_\sys), \state(s_0, \sigma_\sys), \numTicks(\sigma_\sys), D) =
                (\PCom(D_\Cau), s)
        \end{cases}
    \end{align*}
\end{definition}

\begin{definition}
    \label{def:enf-event-system-soundness-property}
    An enforcement event system
    $\EventSystem =
    \left(\Traces_\sys, \Databases \cup \{ \tick, \sskip \}, \RFun, \PFun \right)$
    is \emph{sound} with respect to some
    property $P \subseteq \Traces_{\sys \omega}$ iff
    \begin{equation*}
        \EventSystem, \emptytrace_\sys \vDash P.
    \end{equation*}    
\end{definition}
Equivalently, Definition \ref{def:enf-event-system-soundness-property}
states that an enforcement event system $\EventSystem$ and
$\Initialstates = \{ \emptytrace_\sys \}$
are sound for a property $P \subseteq \Traces_{\sys \omega}$
iff $\EStraces_\omega(\EventSystem, \emptytrace_\sys) \subseteq P$.

\begin{definition}
    \label{def:enf-event-system-soundness-formula}
    An enforcement event system $\EventSystem$ is sound for a formula $\phi$
    iff it is sound for the property $\Language(\phi)$.
\end{definition}

\begin{theorem}
    \label{theo:induced-enforcement-event-system-soundness}
    If an enforcer $\Enforcer = (\EnfStates, s_0, \EnfFun)$ is sound for a
    property $P \subseteq \Traces_\omega$, then the induced enforcement event
    system $\EventSystem(\Enforcer)$ is sound for the property
    $\cConv(P) \subseteq \Traces_{\sys \omega}$.
\end{theorem}

\begin{proof}
    \label{proof:theorem-induced-enforcement-event-system-soundness}
    The proof for Theorem
    \ref{theo:induced-enforcement-event-system-soundness}
    proceeds by structural induction over traces
    $\EStraces_\omega(\EventSystem(\Enforcer), \emptytrace_\sys)$
    of an enforcement event system that is induced by some enforcer
    $\Enforcer = (\EnfStates, s_0, \EnfFun)$ which is sound for
    a property $P \subseteq \Traces_\omega$.
    Specifically we will prove that all
    $\sigma_\sys \in \EStraces(\EventSystem(\Enforcer)), \emptytrace_\sys)$
    are in $\cConv(P) \subseteq \Traces_{\sys \omega}$.

    TODO: argue that for any finite trace of the event system, there exists
    a sequence in c(P) of which the trace is a finite prefix...
\end{proof}


\begin{definition}
    \label{def:enforcement-event-system-transparency}
    TODO: when is an enforcement event system transparent?

    (First finish soundness)
\end{definition}

%%%%%%%%%%%%%%%%%%%%%%%%%%%%%%%%%%%%%%%%%%%%%%%%%%%%%%%%%%%%%%%%%%%%%%%%%%%%%%%%
%                                                                              %
%%%%%%%%%%%%%%%%%%%%%%%%%%%%%%%%%%%%%%%%%%%%%%%%%%%%%%%%%%%%%%%%%%%%%%%%%%%%%%%%
\section{Refinement}
\label{section:refinement}


%%%%%%%%%%%%%%%%%%%%%%%%%%%%%%%%%%%%%%%%%%%%%%%%%%%%%%%%%%%%%%%%%%%%%%%%%%%%%%%%
\subsection{Refinement of Enforcement Event Systems}

For Enforcement Event Systems we will now define refinement
via a mediator function and a simulation relation as in Sprenger et. al \cite{igloo}.

\begin{definition}
    \label{def:enforcement-event-system-refinement}
    An enforcement event system
    $\EventSystem_C = \left( \Traces_{\sys C}, \Databases_C \cup \{\tick, \sskip \}, \RFun_C, \PFun_C \right)$
    for some signature 
    $\Sigma_C = \left( \Domain, \Events_C, \Functions, \iota_C \right)$
    with initial states $\Initialstates_C \subseteq \Statespace_C$
    refines an enforcement event system
    $\EventSystem_A = \left( \Traces_{\sys A}, \Databases_A \cup \{\tick, \sskip \}, \RFun_A, \PFun_A \right)$
    for some signature 
    $\Sigma_A = \left( \Domain, \Events_A, \Functions, \iota_A \right)$
    with initial states $\Initialstates_A \subseteq \Statespace_A$
    via a mediator function
    \begin{equation*}
        \MFun : \Traces_C \times \Eventspace_C \to \Eventspace_A
    \end{equation*}
    iff
    there is a simulation relation
    \begin{equation*}
        \SimRel \subseteq \Statespace_C \times \Statespace_A
    \end{equation*}
    such that
    \begin{enumerate}
        \item
            for all $\sigma_C \in \Initialstates_C$ there exists some
            $\sigma_A \in \Initialstates_A$ such that
            $(\sigma_C, \sigma_A) \in \SimRel$, and
        \item
            for all $\sigma_A \in \Statespace_A$,
            $\sigma_C, \sigma_C' \in \Statespace_C$, and
            $e_C \in \Eventspace_C$
            such that $(\sigma_C, \sigma_A) \in \SimRel$ and
            $\sigma_C \Transitionrelation{}{e_C} \sigma_C'$
            there exists $\sigma_A'$
            such that $(\sigma_C', \sigma_A') \in \SimRel$ and
            $\sigma_A \Transitionrelation{}{\MFun(\sigma_C, e_C)} \sigma_A'$.
    \end{enumerate}
    
\end{definition}
It is important to highlight, that our definition of a mediator function
differs from that in Sprenger et al. \cite{igloo} in that we include a second
argument for the trace up to the current point in time.
The reason for this is how we define MFOTL and Lex refinement in the following
sections.
Without going into the details, on a high level MFOTL refinements can define
let bindings for abstract events, thus "concretizing" them.
And let bindings allow for arbitrary MFOTL formulae on their right hand side which
can of course depend on the past.
That is why we need the history of events as an additional argument for our mediator
function.


%%%%%%%%%%%%%%%%%%%%%%%%%%%%%%%%%%%%%%%%%%%%%%%%%%%%%%%%%%%%%%%%%%%%%%%%%%%%%%%%
\subsection{Refinement of MFOTL Specifications}
\label{subsection:refinement-of-mfotl-specifications}

In this section we define refinement for MFOTL formulas and we proof
that, under our refinement definition,
if a formula $\phi_C$ refines a causable formula $\phi_A : \Cau$,
then $\phi_C$ is enforceable.

\begin{definition}
    \label{def:enforceable-mfotl-specification}
    An \emph{EMFOTL Specification} is a triple
    $S = (\phi, \Sigma, \kappa)$ consisting of an EMFOTL formula $\phi$,
    a signature $\Sigma$, and a capability mapping $\kappa$.
\end{definition}

\begin{definition}
    \label{def:mfotl-refinement}
    Refinement of EMFOTL specifications is defined by the following rules.
    EMFOTL refinement is defined over tuples of specifications and enforcer contexts.
    \begin{gather*}
        \prftree[r]{Ref-Imp}
            {\phi_C \implies \phi_A}
            {\Gamma \vdash \phi_C : \tau_C}
            {\Gamma \vdash \phi_A : \tau_A}
            {\tau_C \capgeq \tau_A}
            % TODO: do the two contexts actually have to be equal, does it matter
            {(\phi_C, \Sigma_C, \kappa_C), \Gamma \refineL (\phi_A, \Sigma_A, \kappa_A), \Gamma} \\
        \prftree[r]{Ref-Conj}
            {S_C, \Gamma_C \refineL (\phi_A, \Sigma_A, \kappa_A), \Gamma_A}
            {S_C, \Gamma_C \refineL (\psi_A, \Sigma_A, \kappa_A), \Gamma_A}
            {S_C, \Gamma_C \refineL (\phi_A \wedge \psi_A, \Sigma_A, \kappa_A), \Gamma_A} \\
        \prftree[r]{Ref-Let}
            {\Gamma_C \vdash \psi_C : \Gamma_A(e)}
            {\forall d \in \Events_C \cap \Events_A .~ \Gamma_C(d) = \Gamma_A(d)}
            {(\LLet{e(\bar{t})}{\psi_C}{\phi_A}, \Sigma_C, \kappa_C), \Gamma_c
                \refineL (\phi_A, \Sigma_A, \kappa_A), \Gamma_A} \\
        \prftree[r]{Ref-Trans}
            {S_C, \Gamma_C \refineL S_B, \Gamma_B}
            {S_B, \Gamma_B \refineL S_A, \Gamma_A}
            {S_C, \Gamma_C \refineL S_A, \Gamma_A}
    \end{gather*}
\end{definition}


\begin{theorem}
    \label{theo:mfotl-refinement-implies-enforcement-event-system-refinement}
    If $S_C, \Gamma_C \refineL S_A, \Gamma_A$ with $S_C = (\phi_C, \Sigma_C, \kappa_C)$ and
    $S_A = (\phi_A, \Sigma_A, \kappa_A)$ and $\phi_A : \Cau$, then
    $\phi_C$ is enforceable.
\end{theorem}
\begin{proof}
    \label{proof:mfotl-refinement-implies-enforcement-event-system-refinement}
    The proof for Theorem \ref{theo:mfotl-refinement-implies-enforcement-event-system-refinement}
    proceeds by structural induction over the refinement rules for MFOTL specifications.
    \begin{itemize}
        \item \textbf{Base Case} (Ref-Imp)

        % Under the Ref-Imp rule, we have that $\phi_C \implies \phi_A$.
        Under the Ref-Imp rule the two enforcer contexts, are identical and the
        the capability of $\phi_C$ under that single context $\Gamma$ is at least as
        powerful as that of $\phi_A$.
        Thus it follows directly that if $\phi_A$ is causable, then $\phi_C$ is at least
        causable and in turn enforceable.
            
        \item \textbf{Induction Hypothesis}

        \item \textbf{Induction Step}

        We make a case distinction on the remaining refinement rules
        \begin{itemize}
            \item \textbf{Ref-Conj}

            Under this rule, $\phi_A = \phi_{A1} \wedge \phi_{A2}$.
            Since $\phi_A$ is causable, each conjunct must be causable individually.
            \item \textbf{Ref-Let}
            \item \textbf{Ref-Trans}
        \end{itemize}
    \end{itemize}
\end{proof}



%%%%%%%%%%%%%%%%%%%%%%%%%%%%%%%%%%%%%%%%%%%%%%%%%%%%%%%%%%%%%%%%%%%%%%%%%%%%%%%%
\subsection{Refinement of CRE Specifications}
\label{subsection:refinement-of-cre-specifications}
In this subsection we recall refinement of Lex specifications as it is described in
Hublet et. al \cite{lex}. We then go ahead and proof that whenever a CRE
specification refines another CRE specification, their compiled MFOTL
formulas are likewise refinements of each other.

\begin{definition}
    \label{def:lex-refinement}
    A specification
    $S_C = \left( \Sigma_C, \Phi_C, \kappa_C \right)$
    with signature $\Sigma_C = \left( \Domain, \Events_C, \Functions, \iota_C \right)$
    and CRE-rules $\Phi_C = \left(C_C, R_C, E_C \right)$
    refines a specification
    $S_A = \left( \Sigma_A, \Phi_A, \kappa_A \right)$
    with signature $\Sigma_A = \left( \Domain, \Events_A, \Functions, \iota_A \right)$
    and CRE-rules $\Phi_A = \left(C_A, R_A, E_A \right)$
    under condition $\Psi = \left(C_C, R_A \setminus R_C, E_C \right)$
    iff all of the following hold
    \begin{enumerate}
        \item 
        $\Events_A \setminus \kappa_A^{-1} 
        (\lbrace\texttt{suppressable},\texttt{observable},\bot\rbrace)
        \subseteq \Events_C$
        and
        for all
        $e \in \Events_A \cap \Events_C,
        \iota_A(e) = \iota_C(e)$
        and
        $\kappa_A(e) \refineL \kappa_C(e)$;
        \item
        $C_A \subseteq C_C$ and
        there exists
        $\Gamma$ such that for all $r:\phi \Rightarrow e(\overline{t}) \in C_C \setminus C_A$,
        all events appearing in $\phi$ are in
        $\Events_C \setminus \Events_A$
        and
        $\Gamma \vdash e(\overline{t}) : \kappa_A(e)$;
        \item
        For all $r: \phi \Rightarrow \exc(s) \in E_A$ such that
        a rule labeled by $s$ exists in
        $\mathcal{S}_C$, $r: \phi \Rightarrow \exc(s) \in E_C$.
    \end{enumerate}
    
\end{definition}

\begin{definition}
    \label{def:lex-inconditional-refinement}
    If we require the following additional condition in Definition \ref{def:lex-refinement}
    \begin{enumerate}
        \setcounter{enumi}{3}
        \item $R_A \subseteq R_C$
    \end{enumerate}
    $S_C$ is an \emph{Incoditional Refinement} of $S_A$.
\end{definition}